\clearpage\pagestyle{otMain}\setlength{\otBlocEspacement}{2.5mm}\fontsize{11}{15}\selectfont\corrigesChapitre{16}{Triangle rectangle et cercle circonscrit}
\begin{blocLivre}{cor-16-1}\begin{corrige}{1}{Trouver le centre}\demarche{Repérer l’hypoténuse, côté opposé à l’angle droit.}
ABC est rectangle en A, donc le centre de son cercle circonscrit est le $\res{\text{milieu de l’hypoténuse }[BC]}$. Le rayon vaut $\frac92=\res{4{,}5\ \mathrm{cm}}$.
\label{fig-cor-16-1}\figureGeo[corrige]{cer-01}
\end{corrige}\end{blocLivre}
\begin{blocLivre}{cor-16-2}\begin{corrige}{2}{Une médiane}\demarche{Utiliser le milieu de l’hypoténuse.}
Dans un triangle rectangle, la médiane issue du sommet de l’angle droit mesure la moitié de l’hypoténuse : $AI=\frac{13}2=\res{6{,}5\ \mathrm{cm}}$.
\label{fig-cor-16-2}\figureGeo[corrige]{cer-02}
\end{corrige}\end{blocLivre}
\begin{blocLivre}{cor-16-3}\begin{corrige}{3}{Un angle droit garanti}\demarche{Vérifier qu’un côté est un diamètre.}
ABC est inscrit dans un cercle dont un diamètre est son côté $[BC]$. D’après la réciproque, $\res{\text{ABC est rectangle en A}}$, sommet opposé au diamètre.
\label{fig-cor-16-3}\figureGeo[corrige]{cer-03}
\end{corrige}\end{blocLivre}
\begin{blocLivre}{cor-16-4}\begin{corrige}{4}{Le rayon par Pythagore}\demarche{Calculer l’hypoténuse, puis sa moitié.}
Le théorème de Pythagore donne $BC=\sqrt{36+64}=10$ cm. Le centre du cercle circonscrit est le milieu de l’hypoténuse : le rayon vaut $\frac{10}2=\res{5\ \mathrm{cm}}$.
\label{fig-cor-16-4}\figureGeo[corrige]{cer-04}
\end{corrige}\end{blocLivre}
\begin{blocLivre}{cor-16-5}\begin{corrige}{5}{Vrai ou faux ?}\demarche{Chercher un contre-exemple.}
Elle est $\res{\text{fausse}}$. Un triangle équilatéral inscrit dans un cercle est un contre-exemple : ses trois angles mesurent $60^\circ$. Aucun de ses côtés n’est un diamètre.
\label{fig-cor-16-5}\figureGeo[corrige]{cer-05}
\end{corrige}\end{blocLivre}
\begin{blocLivre}{cor-16-6}\begin{corrige}{6}{Une médiane qui décide}\demarche{Montrer que A est sur le cercle de diamètre [BC].}
$IA=IB=IC=3{,}5$ cm : A est sur le cercle de centre I et de rayon 3,5 cm, qui a pour diamètre $[BC]$. Le triangle ABC est inscrit dans ce cercle avec $[BC]$ comme diamètre : $\res{\text{il est rectangle en A}}$.
\label{fig-cor-16-6}\figureGeo[corrige]{cer-06}
\end{corrige}\end{blocLivre}
\begin{blocLivre}{cor-16-7}\begin{corrige}{7}{Calculer un angle}\demarche{Trouver d’abord l’angle droit.}
Le triangle ABC est inscrit dans le cercle de diamètre $[AB]$ : il est rectangle en C. La somme des angles donne $\widehat{BAC}=180-90-32=\res{58^\circ}$.
\label{fig-cor-16-7}\figureGeo[corrige]{cer-07}
\end{corrige}\end{blocLivre}
\begin{blocLivre}{cor-16-8}\begin{corrige}{8}{Un rectangle sans équerre}\demarche{Utiliser les propriétés des diagonales.}
Les diagonales $[AC]$ et $[BD]$ ont le même milieu O : ABCD est un parallélogramme. Elles ont la même longueur, celle d’un diamètre, $\res{8\ \mathrm{cm}}$. Un parallélogramme dont les diagonales ont la même longueur est un $\res{\text{rectangle}}$.
\label{fig-cor-16-8}\figureGeo[corrige]{cer-08}
\end{corrige}\end{blocLivre}
\begin{blocLivre}{cor-16-9}\begin{corrige}{9}{Des triangles rectangles en M}\demarche{Utiliser le théorème et sa réciproque.}
Si AMB est rectangle en M, son cercle circonscrit a pour diamètre $[AB]$ : M est sur le cercle de diamètre $[AB]$. Réciproquement, tout point de ce cercle, autre que A et B, convient. Ce sont les points du $\res{\text{cercle de centre le milieu de }[AB]\text{ et de rayon 3 cm}}$, sauf A et B.
\label{fig-cor-16-9}\figureGeo[corrige]{cer-09}
\end{corrige}\end{blocLivre}
\begin{blocLivre}{cor-16-10}\begin{corrige}{10}{Avec Pythagore}\demarche{Identifier l’angle droit avant d’appliquer Pythagore.}
ABC est inscrit dans le cercle de diamètre $[AB]$ : il est rectangle en C, avec $AB=10$ cm. Pythagore donne $BC^2=10^2-6^2=64$, donc $\res{BC=8\ \mathrm{cm}}$.
\label{fig-cor-16-10}\figureGeo[corrige]{cer-10}
\end{corrige}\end{blocLivre}
\begin{blocLivre}{cor-16-11}\begin{corrige}{11}{Une perpendiculaire au compas}\demarche{Reconnaître un triangle inscrit dans un demi-cercle.}
A est sur le cercle de diamètre $[BC]$, distinct de B et C : le triangle ABC est rectangle en A. Donc $(AC)\perp(AB)$, c’est-à-dire $\res{(AC)\perp(d)}$. La construction n’utilise que la règle et le compas.
\label{fig-cor-16-11}\figureGeo[corrige]{cer-11}
\end{corrige}\end{blocLivre}
\begin{blocLivre}{cor-16-12}\begin{corrige}{12}{Dans un demi-cercle}\demarche{Combiner l’angle droit et les côtés égaux.}
ABC est rectangle en C et isocèle. Pythagore donne $AC^2+BC^2=100$, donc $2AC^2=100$ et $AC^2=50$ : $\res{AC=\sqrt{50}\ \mathrm{cm}}\approx\res{7{,}07\ \mathrm{cm}}$. L’aire vaut $\frac{AC\times BC}2=\frac{50}2=\res{25\ \mathrm{cm}^2}$.
\label{fig-cor-16-12}\figureGeo[corrige]{cer-12}
\end{corrige}\end{blocLivre}
\begin{blocLivre}{cor-16-13}\begin{corrige}{13}{Distance du centre à un côté}\demarche{Utiliser le cercle puis la droite des milieux.}
$BC=\sqrt{25+144}=13$ cm, donc le rayon vaut $\res{6{,}5\ \mathrm{cm}}$, et O est le milieu de $[BC]$. Dans ABC, O et J sont les milieux de $[BC]$ et $[AB]$ : $OJ=\frac{AC}2=\res{6\ \mathrm{cm}}$ et $(OJ)\parallel(AC)$, donc $(OJ)\perp(AB)$.
\label{fig-cor-16-13}\figureGeo[corrige]{cer-13}
\end{corrige}\end{blocLivre}
\begin{blocLivre}{cor-16-14}\begin{corrige}{14}{Un rectangle inscrit}\demarche{Reconnaître que la diagonale est un diamètre.}
Les diagonales d’un rectangle ont le même milieu et la même longueur : ce milieu est le centre du cercle et $AC=13$ cm est un diamètre. ABC est rectangle en B, donc $BC=\sqrt{13^2-5^2}=\res{12\ \mathrm{cm}}$. L’aire vaut $5\times12=\res{60\ \mathrm{cm}^2}$.
\label{fig-cor-16-14}\figureGeo[corrige]{cer-14}
\end{corrige}\end{blocLivre}
\begin{blocLivre}{cor-16-15}\begin{corrige}{15}{Sans mesurer d’angle}\demarche{Démontrer l’angle droit, puis appliquer Pythagore.}
O est sur $[BC]$ et $OB=OC=4$ cm : $[BC]$ est un diamètre de 8 cm. Le triangle ABC, inscrit dans le cercle de diamètre $[BC]$, est rectangle en A. Donc $AC^2=64-25=39$ : $\res{AC=\sqrt{39}\ \mathrm{cm}}\approx\res{6{,}24\ \mathrm{cm}}$. L’angle droit est démontré par la réciproque, pas mesuré.
\label{fig-cor-16-15}\figureGeo[corrige]{cer-15}
\end{corrige}\end{blocLivre}
